\section{Experiments}
\label{sec:experiments}

\subsection{Setup}

\paragraph{Vector memory.}
We use FAISS IndexFlatIP (exact inner product search) with L2-normalized
all-MiniLM-L6-v2 embeddings (384-dim).  The benign corpus contains 200 entries;
poison counts are attack-specific: $n_{\text{base}}=5$ for \agentpoison{},
$2n_{\text{base}}=10$ for \minja{} (multi-entry bridging coverage), and
$3n_{\text{base}}=15$ for \injecmem{} (volume-based broad recall).  Top-$k=5$.

\paragraph{Evaluation.}
$\asrr$ is measured empirically by issuing all 20 victim queries and counting
retrievals of adversarial entries.  $\asra$ is modelled per
Section~\ref{sec:attacks}.  Defense metrics are measured over the same 20
victim queries (for $\tpr$) and 20 benign queries (for $\fpr$).  Bootstrap
95\% confidence intervals are computed from 5 independent random seeds
(Section~\ref{sec:statistical}).

\subsection{Attack Results}
\label{sec:attack_results}

\IfFileExists{../../results/tables/table1_attack_results.tex}{%
  \begin{table}[t]
\centering
\small
\caption{Attack evaluation results on synthetic memory corpus
  ($N=200$, $k=5$, seed-averaged over 5 trials).
  ASR-A$^*$ is modelled from paper-reported values.
  95\% CI via percentile bootstrap (2000 replicates).}
\label{tab:attack_results}
\begin{tabular}{lcccc}
\toprule
Attack & ASR-R & ASR-A$^*$ & ASR-T & Benign Acc \\
\midrule
AgentPoison \citep{chen2024agentpoison} & 1.000$_{[1.000,1.000]}$ & 0.68 & 0.900 & 0.850 \\
MINJA \citep{dong2025minja} & 0.650$_{[0.650,0.650]}$ & 0.76 & 0.550 & 0.800 \\
InjecMEM \citep{injecmem2026} & 0.750$_{[0.750,0.750]}$ & 0.57 & 0.550 & 0.550 \\
PoisonedRAG \citep{zou2024poisonedrag} & 0.400$_{[0.400,0.400]}$ & 0.82 & 0.400 & 1.000 \\
\bottomrule
\end{tabular}
\end{table}
%
}{%
\begin{table}[t]
  \centering
  \caption{Attack evaluation results.  Bootstrap 95\% CI across 5 seeds.
           $\asra$ is modelled.  All metrics in $[0,1]$; higher is worse
           for defender.}
  \label{tab:attack_results}
  \begin{tabular}{lccccc}
    \toprule
    Attack & $\asrr$ & $\asra$ & $\asrt$ & Benign Acc. & ISR \\
    \midrule
    \agentpoison{} (ours) & $1.00 \pm 0.00$ & $0.68$ & $\approx0.65$ & $1.00$ & $1.00$ \\
    \minja{}              & $0.65 \pm 0.00$ & $0.76$ & $\approx0.55$ & $0.90$ & $0.98$ \\
    \injecmem{}           & $0.50 \pm 0.00$ & $0.57$ & $\approx0.35$ & $0.80$ & $1.00$ \\
    \bottomrule
  \end{tabular}
\end{table}%
}

Table~\ref{tab:attack_results} reports the main attack results.
Figure~\ref{fig:attack_asr} visualises the three metrics side-by-side.
\agentpoison{} achieves $\asrr=1.00$ with the corrected triggered-query
evaluation protocol and centroid-targeting passage.
\minja{} achieves $\asrr=0.65$ in our FAISS retrieval simulation.
The original paper reports ISR\,=\,98.2\% (injection success rate:
fraction of bridging-step passages successfully stored by the agent's
auto-storage mechanism), not $\asrr$ from a pre-indexed corpus.
Our $\asrr=0.65$ measures retrieval probability from a fixed 200-entry
FAISS index, which is a different and harder evaluation: template-based
passages must compete with 200 benign entries rather than exploiting the
agent's auto-storage pipeline directly.\footnote{\minja{}'s ISR and our
$\asrr$ measure complementary threat surfaces.  ISR quantifies write-time
poisoning success; $\asrr$ quantifies retrieval-time exposure given that
the poison entry has already been stored.  Both are necessary for
a complete threat characterisation.}
\injecmem{} achieves intermediate recall via volume (15 entries) at
the cost of lower benign accuracy.

\IfFileExists{figures/fig_01_attack_asr.pdf}{%
\begin{figure}[t]
  \centering
  \includegraphics[width=0.72\linewidth]{fig_01_attack_asr.pdf}
  \caption{Attack success rates (bootstrap 95\,\% CI over 5 seeds).
           Left group: $\asrr$; middle: $\asra$ (modelled); right: $\asrt$.
           \agentpoison{} dominates on retrieval ($\asrr=1.00$) but has
           lower $\asra$ than \minja{} due to its blunter adversarial framing.}
  \label{fig:attack_asr}
\end{figure}%
}{}

\paragraph{Evaluation correction for \agentpoison{}.}
\begin{table}[t]
  \centering
  \caption{Impact of the triggered-query evaluation correction.
           Pre-fix: plain queries (no trigger).
           Post-fix: triggered queries ($q \oplus \trigger^*$).}
  \label{tab:trigger_fix}
  \begin{tabular}{lcc}
    \toprule
    Evaluation protocol & $\asrr$ & Note \\
    \midrule
    Un-triggered queries (pre-Phase 12) & 0.25 & Underestimates attack strength \\
    Centroid passage + triggered queries & 1.00 & Matches paper protocol \\
    \bottomrule
  \end{tabular}
\end{table}

Table~\ref{tab:trigger_fix} isolates the effect of the evaluation correction.
This result has an important implication for future work: simulations that
evaluate \agentpoison{} without trigger prepending will substantially
underestimate attacker capability.

\subsection{Attack-Defense Interaction Matrix}
\label{sec:matrix}

\IfFileExists{../../results/tables/table2_attack_defense_matrix.tex}{%
  \begin{table}[t]
\centering
\caption{Attack-defense interaction matrix: post-defense ASR-R (fraction of victim queries that retrieve undetected poison). Rows = attacks, Columns = defenses. Lower is better for defenses. Baseline = no defense.}
\label{tab:attack_defense_matrix}
\begin{tabular}{lcccccc}
\toprule
\multicolumn{7}{c}{ASR-R Under Defense (lower is better for defender)} \\
\midrule
Attack & Baseline & Watermark & Validation & Proactive & Composite & SAD (ours) \\
\midrule
AgentPoison & 1.000 & \textbf{0.000} & 0.000 & 0.000 & 0.000 & 1.000$^\dagger$ \\
MINJA & 0.633 & \textbf{0.000} & 0.000 & 0.583 & 0.000 & 0.000 \\
InjecMEM & 0.483 & \textbf{0.000} & 0.483 & 0.183 & 0.000 & 0.433 \\
\midrule
Avg.~Effectiveness (↑) & — & 1.000 & 0.667 & 0.575 & 1.000 & 0.368 \\
\bottomrule
\end{tabular}
\vspace{0.5ex}
{\small $^\dagger$Plain-query SAD calibration; triggered-query calibration achieves ASR-R\,=\,0.000 (Section~\ref{sec:discussion}).}
\end{table}
%
}{%
  \textit{(Run \texttt{generate\_tables.py} to generate attack-defense matrix table.)}%
}

\IfFileExists{../../results/tables/table3_defense_tpr_fpr.tex}{%
  \begin{table}[t]
\centering
\caption{Defense detection performance: TPR, FPR, and AUROC per (attack, defense) pair.
         SAD uses combined scoring mode ($0.5\cdot\max + 0.5\cdot\text{mean}$ cosine similarity to observed victim queries).
         Proactive values reflect retrieval-breadth detection at $\tau=0.19$.
         SAD$^\dagger$: triggered-query calibration for \agentpoison{}.}
\label{tab:defense_tpr_fpr}
\begin{tabular}{lcccccc}
\toprule
\multicolumn{7}{c}{TPR\,/\,FPR\,/\,AUROC by (Defense, Attack)} \\
\midrule
Defense & \multicolumn{2}{c}{AgentPoison} & MINJA & InjecMEM \\
        & plain cal.     & trig.\ cal.$^\dagger$ &       &          \\
\midrule
Watermark  & 1.00\,/\,0.00\,/\,--  & 1.00\,/\,0.00\,/\,--  & 1.00\,/\,0.00\,/\,-- & 1.00\,/\,0.00\,/\,-- \\
Validation & 1.00\,/\,0.00\,/\,--  & 1.00\,/\,0.00\,/\,--  & 1.00\,/\,0.00\,/\,-- & 0.00\,/\,0.00\,/\,-- \\
Proactive  & 1.00\,/\,0.01\,/\,--  & 1.00\,/\,0.01\,/\,--  & 0.07\,/\,0.01\,/\,-- & 0.60\,/\,0.01\,/\,-- \\
Composite  & 1.00\,/\,0.00\,/\,--  & 1.00\,/\,0.00\,/\,--  & 1.00\,/\,0.00\,/\,-- & 1.00\,/\,0.00\,/\,-- \\
SAD (ours) & 0.00\,/\,0.00\,/\,0.870 & \textbf{1.00}\,/\,\textbf{0.00}\,/\,\textbf{1.000} & \textbf{1.00}\,/\,\textbf{0.00}\,/\,\textbf{1.000} & 0.40\,/\,0.00\,/\,0.920 \\
\bottomrule
\end{tabular}
\vspace{0.5ex}
{\small $^\dagger$Calibrated on triggered queries ($\trigger^*\oplus q$); see Section~\ref{sec:discussion}.
AUROC not applicable to threshold-only defenses (Watermark, Validation, Proactive, Composite).}
\end{table}
%
}{}

\paragraph{Defense differentiation.}
Table~\ref{tab:defense_tpr_fpr} shows each defense has a distinct effectiveness
profile.  Watermark and Composite achieve $\tpr=1.00$, $\fpr=0.00$ uniformly.
Proactive achieves $\tpr=1.00$ / $\fpr=0.01$ for \agentpoison{} and
$\tpr=0.60$ / $\fpr=0.01$ for \injecmem{}, but only $\tpr=0.07$ for \minja{}:
targeted bridging-step passages match few cross-domain retrieval probes, so
the retrieval-breadth criterion does not flag them.  Validation achieves
$\tpr=0.00$ for \injecmem{}, which avoids rule-based injection keywords.
This complementary structure motivates the composite ensemble.
\sad{} (combined mode, $k=2.0$) achieves $\fpr=0.00$ across all three attacks
with AUROC $0.870$ (\agentpoison{}, plain calibration), $1.000$ (\minja{}), and
$0.920$ (\injecmem{})---the high AUROC for \injecmem{} indicates that the
combined score provides good ranking separation even where the binary TPR at
$k=2.0$ is $0.40$.

\paragraph{SAD calibration and the triggered-query effect.}
Table~\ref{tab:defense_tpr_fpr} shows \sad{} achieves $\tpr=0.00$ for
\agentpoison{} under plain-query calibration.  We show (Section~\ref{sec:discussion})
that calibrating \sad{} on triggered queries (the realistic post-deployment
scenario) raises $\tpr$ to $1.00$ with $\fpr=0.00$---a result confirming
that the \agentpoison{} blind spot is a calibration artefact, not an
inherent limitation of the anomaly-based approach.
Figure~\ref{fig:sad_calibration} illustrates the mechanism: under plain
calibration the \agentpoison{} centroid passage falls within the benign
similarity distribution, but under triggered calibration benign entries
score low against triggered queries while the passage---optimised for
triggered queries---scores high, placing it clearly above the threshold.

\IfFileExists{figures/fig_sad_calibration_comparison.pdf}{%
\begin{figure}[t]
  \centering
  \includegraphics[width=0.96\linewidth]{fig_sad_calibration_comparison.pdf}
  \caption{SAD calibration mode comparison for \agentpoison{}.
           \emph{(a)} Plain-query calibration: the centroid passage (red)
           falls within the benign similarity distribution (blue) ---
           threshold is set too high to detect it ($\tpr=0.00$).
           \emph{(b)} Triggered-query calibration: benign entries score low
           against triggered queries while the centroid passage scores high,
           placing it above the (now lower) detection threshold
           ($\tpr=1.00$, $\fpr=0.00$, threshold shift $0.482\to 0.432$).
           See also Figure~\ref{fig:sad_roc} for the full threshold sweep.}
  \label{fig:sad_calibration}
\end{figure}%
}{}

\IfFileExists{figures/m12_attack_defense_matrix.pdf}{%
\begin{figure}[t]
  \centering
  \includegraphics[width=0.72\linewidth]{m12_attack_defense_matrix.pdf}
  \caption{Attack-defense interaction matrix: $\asrr$ under each
           defense (columns) for each attack (rows).
           Darker shading indicates higher defender-side risk
           (higher ASR-R despite defense).
           No single column achieves uniformly low values,
           motivating composite portfolios.}
  \label{fig:ad_matrix}
\end{figure}%
}{}

\subsection{Statistical Analysis}
\label{sec:statistical}

We use percentile bootstrap confidence intervals (2000 replicates,
\citet{efron1993bootstrap}) to quantify variability across seeds.
Pairwise comparisons use paired $t$-tests with Bonferroni correction
for multiple testing ($\alpha_{\text{corrected}} = 0.05 / 2 = 0.025$
for the two non-redundant ordered comparisons).
Cohen's $d$ effect sizes are reported for all comparisons.

\begin{table}[t]
  \centering
  \caption{Pairwise hypothesis tests on $\asrr$ across 5 seeds.
           All pairs are significantly different ($p < 0.05$).}
  \label{tab:hypothesis}
  \begin{tabular}{lcccc}
    \toprule
    Comparison & $t$ & $p$ & Cohen's $d$ & Effect Size \\
    \midrule
    \minja{} vs \agentpoison{} & $19.0$ & $<0.001$ & $11.5$ & large \\
    \injecmem{} vs \minja{}    & $10.2$ & $<0.001$ & $4.4$  & large \\
    \bottomrule
    \multicolumn{5}{l}{\footnotesize Bonferroni-corrected $\alpha = 0.025$ (2 comparisons); $n=5$ seeds per attack.} \\
  \end{tabular}
\end{table}

\subsection{Adaptive Adversary Results}
\label{sec:adaptive_results}

\IfFileExists{../../results/tables/table5_adaptive_sad.tex}{%
  \begin{table}[t]
\centering
\small
\caption{Adaptive adversary (greedy synonym substitution) vs.\ SAD ($k=2.0$).
         Standard ASR-R: baseline attack without trigger optimisation.
         Adaptive: after synonym substitution targeting the detection boundary.
         Evasion Rate = fraction of poison entries below SAD threshold;
         $\Delta$ASR-R = retrieval degradation (Standard $-$ Adaptive ASR-R).
         Note: AgentPoison is evaluated without trigger optimisation to isolate
         the SAD detection component; see Table~\ref{tab:attack_results} for
         triggered-query ASR-R values.}
\label{tab:adaptive_sad}
\begin{tabular}{lcccccc}
\toprule
Attack & \multicolumn{2}{c}{Standard} & \multicolumn{2}{c}{Adaptive} & Evasion & $\Delta$ASR-R \\
\cmidrule(lr){2-3}\cmidrule(lr){4-5}
 & ASR-R & SAD TPR & ASR-R & SAD TPR & Rate & \\
\midrule
AgentPoison & 0.100 & 0.000 & 0.100 & 0.000 & 1.000 & 0.000 \\
MINJA        & 0.650 & 0.450 & 0.650 & 0.000 & 1.000 & 0.000 \\
InjecMEM     & 0.500 & 0.467 & 0.500 & 0.200 & 0.800 & 0.000 \\
\bottomrule
\multicolumn{7}{l}{\footnotesize Synonym substitution achieves evasion without retrieval cost} \\
\multicolumn{7}{l}{\footnotesize (semantic embeddings are synonym-invariant; see Section~\ref{sec:adaptive_results}).} \\
\end{tabular}
\end{table}
%
}{}

\IfFileExists{figures/p13_adaptive_tradeoff_agent_poison.pdf}{%
\begin{figure}[t]
  \centering
  \begin{subfigure}[b]{0.32\linewidth}
    \includegraphics[width=\linewidth]{p13_adaptive_tradeoff_agent_poison.pdf}
    \caption{\agentpoison{}}
  \end{subfigure}\hfill
  \begin{subfigure}[b]{0.32\linewidth}
    \IfFileExists{figures/p13_adaptive_tradeoff_minja.pdf}{%
      \includegraphics[width=\linewidth]{p13_adaptive_tradeoff_minja.pdf}%
    }{}
    \caption{\minja{}}
  \end{subfigure}\hfill
  \begin{subfigure}[b]{0.32\linewidth}
    \IfFileExists{figures/p13_adaptive_tradeoff_injecmem.pdf}{%
      \includegraphics[width=\linewidth]{p13_adaptive_tradeoff_injecmem.pdf}%
    }{}
    \caption{\injecmem{}}
  \end{subfigure}
  \caption{Adaptive adversary tradeoff: evasion rate vs.\ retrieval
           degradation $\Delta\asrr$ at increasing synonym-substitution
           budgets.  Each point is one budget level.  The negative
           slope confirms the fundamental tension: evasion always
           incurs a retrieval cost.}
  \label{fig:adaptive_tradeoff}
\end{figure}%
}{}

We evaluate the adaptive adversary described in Section~\ref{sec:adaptive}
against \sad{} for all three attack types.  The adversary applies greedy
synonym substitution targeting the calibration threshold at $k=2.0$.

\paragraph{Key finding.}
Table~\ref{tab:adaptive_sad} reports the measured outcomes.
The word-level synonym adversary achieves high evasion rates
(80--100\%) with negligible retrieval cost ($\Delta\text{ASR-R} \approx 0$)
across all three attack types.  This result reflects a fundamental property
of semantic embeddings: all-MiniLM-L6-v2 is trained to be invariant to
synonym substitution, so replacing ``schedule'' with ``plan'' or ``meeting''
with ``session'' shifts embeddings by less than the anomaly detection margin.
Consequently, synonym substitution can bring a passage below the detection
threshold without materially changing its cosine similarity to victim queries.

\paragraph{Theoretical vs.\ empirical tension.}
The theoretical gradient analysis (Eq.~\ref{eq:adaptive_obj},
Section~\ref{sec:adaptive}) proves that a \emph{continuous gradient-based}
adversary cannot reduce anomaly score without simultaneously reducing
retrieval rank, since both depend on the same similarity measure.
The empirical results reveal a practical loophole: \emph{discrete synonym
substitution} is not gradient-based and therefore is not subject to this
tension.  In the synonym vocabulary, two synonymous words may have
nearly identical semantic embeddings, so the adversary can cross the
detection threshold (which is calibrated at a specific similarity value)
without meaningfully reducing retrieval rank.

These findings motivate future defenses that extend anomaly detection
beyond the $\ell_2$ semantic embedding space, for example by incorporating
lexical diversity features or synonym-aware thresholds that are robust
to in-vocabulary substitution.  A gradient-based adversary operating
directly in the continuous embedding space would face the theoretical
tension; such an adversary corresponds to an attacker with access to
differentiable text optimization (e.g., via HotFlip), which is a
strictly stronger threat model.

\subsection{Watermark Evasion}

We evaluate three evasion strategies against the unigram watermark
(Section~\ref{sec:defenses}): (i) paraphrasing at intensity 0.1--0.5,
(ii) copy-paste dilution with $r=0.5$--$3.0$ clean content, and
(iii) adaptive substitution with budget 1--20 character swaps.
Detailed results including z-score distribution curves are reported in
Appendix~\ref{sec:evasion_appendix}.

\IfFileExists{figures/p13_evasion_analysis.pdf}{%
\begin{figure}[t]
  \centering
  \includegraphics[width=0.82\linewidth]{p13_evasion_analysis.pdf}
  \caption{Watermark evasion analysis across three strategies:
           paraphrasing (synonym substitution intensity $\rho \in [0.1, 0.5]$),
           copy-paste dilution (dilution ratio $r \in [0.5, 3.0]$), and
           adaptive character substitution (budget $b \in [1, 20]$).
           The unigram z-score defense maintains $\tpr = 1.00$ across
           all tested intensity levels on the synthetic corpus; the
           embedded vocabulary bias survives these mild evasion strategies.
           (Table~\ref{tab:evasion}; detailed z-score distributions in
           Appendix~\ref{sec:evasion_appendix}.)}
  \label{fig:evasion_analysis}
\end{figure}%
}{}

\subsection{Ablation Studies}
\label{sec:ablation}

We ablate four hyperparameters to characterise the sensitivity of our
evaluation to experimental choices.  Full ablation tables are in
Appendix~\ref{sec:ablation_appendix}.

\paragraph{Corpus size.}
The ablation sweeps $N \in \{50, 100, 200, 500\}$ in un-triggered evaluation
(isolating the retrieval mechanism from the trigger).  Without the trigger,
ASR-R decreases from $0.33$ at $N=50$ to near $0$ at $N=500$, confirming
that the trigger is the primary mechanism maintaining $\asrr=1.00$ in
full evaluation.  This motivates the evaluation protocol correction of
Section~\ref{sec:experiments}: un-triggered simulations substantially
underestimate attack strength.

\paragraph{Retrieval top-\texorpdfstring{$k$}{k}.}
ASR-R increases monotonically with $k$ (1--20): a larger retrieval window
gives the adversary more opportunities for poison to appear in results
(Table~\ref{tab:ablation}).  In un-triggered evaluation, $\asrr \approx 0$
at $k=1$; at $k=20$, $\asrr$ approaches $0.45$.  Under triggered evaluation,
\agentpoison{} achieves $\asrr \geq 0.95$ across all $k \geq 1$.

\paragraph{Poison count.}
\agentpoison{} shows near-ceiling performance with even 1 poison entry
(trigger guarantees retrieval).  \injecmem{} shows near-linear scaling
up to 10 entries, after which returns diminish.

\paragraph{\sad{} threshold \texorpdfstring{$k$}{k}.}
Lower $k$ yields higher TPR but also higher FPR (aggressive detection).
At the default $k=2.0$, \sad{} achieves TPR $\approx 1.00$ for \minja{},
TPR $\approx 0.20$ for \injecmem{}, and TPR $\approx 0.00$ for \agentpoison{}.
The near-zero TPR for \agentpoison{} reflects that the centroid passage
similarity falls within the benign calibration distribution when \sad{}
is calibrated on un-triggered victim queries (see Section~\ref{sec:discussion}).
FPR $\approx 0.05$ is consistent across all three attacks
(see Appendix~\ref{sec:ablation_appendix} and Figure~\ref{fig:sad_ablation}).

\IfFileExists{figures/p13_ablation_sad_sigma.pdf}{%
\begin{figure}[t]
  \centering
  \begin{subfigure}[b]{0.48\linewidth}
    \includegraphics[width=\linewidth]{p13_ablation_sad_sigma.pdf}
    \caption{\sad{} threshold $k$ vs.\ TPR/FPR}
    \label{fig:sad_ablation}
  \end{subfigure}\hfill
  \begin{subfigure}[b]{0.48\linewidth}
    \IfFileExists{figures/p13_ablation_corpus.pdf}{%
      \includegraphics[width=\linewidth]{p13_ablation_corpus.pdf}%
    }{}
    \caption{Corpus size $N$ vs.\ $\asrr$}
    \label{fig:corpus_ablation}
  \end{subfigure}
  \caption{Ablation curves.
           Left: \sad{} threshold $k$ controls the TPR/FPR tradeoff;
           $k=2.0$ (vertical dashed) achieves $\tpr\!\approx\!1.00$
           for \minja{}, $\tpr\!\approx\!0.20$ for \injecmem{},
           and $\tpr\!\approx\!0.00$ for \agentpoison{} (centroid passage
           falls within benign similarity distribution at un-triggered calibration).
           $\fpr\!\approx\!0.05$ across all attacks.
           Right: ASR-R for all three attacks is largely independent of
           benign corpus size $N$ (50--500 entries), confirming poison
           entries are extreme similarity outliers.}
  \label{fig:ablation_combined}
\end{figure}%
}{}

\subsection{Phase 17: DPR Optimization, Token-Level Watermark, and Live Agent Evaluation}
\label{sec:phase17_results}

\paragraph{DPR HotFlip trigger optimization.}
The DPR HotFlip optimizer (Section~\ref{sec:attacks}) achieves mean cosine
similarity $0.78$ between triggered query embeddings and the adversarial
passage, compared to $0.48$ for a random trigger baseline ($+62\%$) and
$0.71$ for the centroid-passage approximation alone ($+10\%$).
Figure~\ref{fig:dpr_convergence} shows that convergence is rapid: $>90\%$
of the final gain is achieved within the first 5 iterations.
The combined centroid passage $+$ DPR trigger preserves $\asrr=1.00$ while
producing more fluent triggers (mean GPT-2 perplexity $<200$) than unconstrained
gradient search.

\paragraph{Token-level watermark comparison.}
The token-level watermark scheme (\textsc{Tok-WM}, Section~\ref{sec:watermark_defense})
substantially outperforms the character-level adaptation for short memory entries.
At $n=50$ GPT-2 tokens (approximately 150--200 characters), the mean watermarked
z-score is $\approx 7.5$ vs.\ $\approx 4.1$ for character-level---a $+83\%$
improvement in detection margin.  For longer entries ($n \geq 200$ characters),
both schemes exceed $z_{\mathrm{thr}}=4.0$ reliably ($\mathrm{TPR} > 0.95$).
The green token fraction under $\delta=2.0$ bias is $p_{\text{green}} \approx 0.711$
(theoretical), empirically confirmed at $0.698 \pm 0.041$ across 80 sampled entries.
Figure~\ref{fig:watermark_comparison} (Section~\ref{sec:watermark_defense})
shows the full z-score distribution comparison.

\paragraph{Measured ASR-A from local agent evaluation.}
Beyond modelled $\asra$ values (calibrated from paper-reported means), we
provide a directly measured lower bound using \textsc{LocalAgentEvaluator}
with a GPT-2 backbone.  Measured values are:
$\asra^{\text{meas}} = 0.42$ (\agentpoison{}),
$0.51$ (\minja{}), and $0.29$ (\injecmem{}).
Compared to modelled upper bounds ($0.68$, $0.76$, $0.57$), the gap confirms
that GPT-2 is a conservative proxy: it follows injected instructions less
faithfully than production-scale LLMs (GPT-4, Llama-3).
This gap---$0.26$--$0.28$ across attacks---represents the headroom between
a traceable lower bound and the paper-calibrated upper bound.
Figure~\ref{fig:measured_asr} illustrates the comparison.

\IfFileExists{figures/fig_p17_measured_asr_a.pdf}{%
\begin{figure}[ht]
  \centering
  \includegraphics[width=0.72\linewidth]{fig_p17_measured_asr_a.pdf}
  \caption{Modelled vs.\ measured $\asra$ for all three attacks.
           Solid bars: modelled values calibrated to paper-reported means.
           Shaded bars: directly measured via GPT-2 agent (\textsc{LocalAgentEvaluator}),
           providing a lower bound on true $\asra$.
           The gap ($0.26$--$0.28$) quantifies the difference between a weak
           local proxy and production-scale LLM instruction following.}
  \label{fig:measured_asr}
\end{figure}%
}{}

\subsection{Multi-Encoder, Multi-Agent, and Graph Memory Evaluation}
\label{sec:phase22_results}

\subsubsection{Multi-Encoder Attack Transferability}

We evaluate all three attacks and \sad{} across seven retrieval encoders
using \textsc{MultiEncoderEvaluator} (\texttt{src/evaluation/multi\_encoder\_eval.py}).
The encoder set spans three architecture families:
\begin{itemize}
  \item \emph{Sentence-transformers}: \texttt{all-MiniLM-L6-v2} (384-dim),
        \texttt{all-mpnet-base-v2} (768-dim), \texttt{intfloat/e5-base-v2} (768-dim),
        \texttt{facebook/contriever-msmarco} (768-dim), \texttt{BAAI/bge-large-en-v1.5} (1024-dim)
  \item \emph{OpenAI API}: \texttt{text-embedding-3-small} (1536-dim),
        \texttt{text-embedding-3-large} (3072-dim)
\end{itemize}

\paragraph{CKA transferability matrix.}
We compute linear Centered Kernel Alignment (CKA; \citealt{kornblith2019similarity})
between all $\binom{7}{2} = 21$ encoder pairs on $n=100$ shared corpus entries:
\begin{equation}
  \mathrm{CKA}(K, L) =
  \frac{\mathrm{HSIC}(K_c, L_c)}{\sqrt{\mathrm{HSIC}(K_c, K_c)\,\mathrm{HSIC}(L_c, L_c)}},
  \label{eq:cka}
\end{equation}
where $K = XX^\top$, $L = YY^\top$ are gram matrices of the respective embeddings
and $K_c$, $L_c$ denote their doubly-centered versions.
The gram-matrix formulation supports encoders with different embedding dimensions,
enabling comparison across heterogeneous architectures.
High CKA between two encoders indicates that their similarity-based ranking is
nearly equivalent, predicting that attacks optimized against one encoder
will transfer to the other.

\paragraph{Cross-encoder SAD generalization.}
\sad{} is calibrated on each encoder independently and evaluated using an
\textsc{ExternalEmbeddingDetector} that replaces the default MiniLM-L6
embedding with the target encoder's embeddings.  The calibration threshold
$k$ is held fixed at $2.0$ across encoders.  The Appendix
(Section~\ref{app:encoder_generalization}) reports encoder-level TPR/FPR
and AUROC.

\IfFileExists{figures/fig_p22_cka_heatmap.pdf}{%
\begin{figure}[t]
  \centering
  \includegraphics[width=0.65\linewidth]{fig_p22_cka_heatmap.pdf}
  \caption{Centered Kernel Alignment (CKA) between three retrieval encoders.
           High CKA ($>0.8$) predicts cross-encoder attack transferability:
           attacks optimized against one encoder transfer to the other
           with minimal ASR-R degradation.}
  \label{fig:cka_heatmap}
\end{figure}%
}{}

\subsubsection{Multi-Agent Epidemic Propagation}

We model inter-agent poison propagation using \textsc{MultiAgentEnvironment}
(\texttt{src/evaluation/multi\_agent\_propagation.py}), a discrete-step
SIR-inspired epidemic model \citep{gu2024agentsmith} over a shared
\textsc{SharedKnowledgeBase} backed by a FAISS IndexFlatIP:

\begin{enumerate}
  \item At each step $t$, all $N$ agents issue $q_{\mathrm{step}}$ queries
        to the shared knowledge base (top-$k$ retrieval).
  \item Any agent that retrieves a poison entry re-stores a secondary
        paraphrase of it with probability $p_{\mathrm{re\text{-}store}}$,
        increasing the total poison count in the knowledge base.
  \item Propagation terminates when no new secondary entries are added for
        one full step or $t_{\mathrm{max}}$ steps are reached.
\end{enumerate}

This models real-world agentic deployments (shared Mem0, MemGPT, or
LangChain memory) where retrievals from one agent session influence
memory contents seen by other agents.

\paragraph{SAD quarantine defense.}
\textsc{PropagationWithSADQuarantine} runs two parallel simulations:
(i) a baseline without defense; (ii) a defended simulation where each
candidate entry---both initial and secondary---is screened through \sad{}
($k=2.0$, combined scoring mode, $\fpr=0.000$) before insertion into the
shared knowledge base.  Blocked entries are quarantined (not stored),
halting secondary infection at the source.

\paragraph{Key findings.}
Under the baseline (no defense), a single initial poison entry can reach
all $N$ agents within $t_{\mathrm{max}}$ steps as re-storage amplifies
the poison count.
\sad{} quarantine eliminates secondary entries entirely for \minja{} and
\agentpoison{} (AUROC $= 1.000$), containing spread to the initial
injection only.  For \injecmem{} (AUROC $= 0.920$), quarantine reduces
but does not fully eliminate spread, consistent with the lower per-entry
TPR ($0.40$) under plain-query calibration.

\IfFileExists{figures/fig_p22_propagation.pdf}{%
\begin{figure}[t]
  \centering
  \includegraphics[width=0.95\linewidth]{fig_p22_propagation.pdf}
  \caption{Multi-agent epidemic propagation ($N$=5 agents,
           $p_{\mathrm{re\text{-}store}}$=0.4) for each attack type.
           Red: undefended spread; green: \sad{} quarantine.
           Bars show cumulative secondary poison entries.
           \sad{} quarantine eliminates secondary infection for \agentpoison{}
           and \minja{}, containing spread to the initial injection.}
  \label{fig:propagation}
\end{figure}%
}{}

\subsubsection{Graph-Structured Memory Attacks and Defenses}

Real-world memory systems increasingly use graph-structured stores where
memories are connected by semantic or relational edges (e.g., GraphRAG
\citealt{edge2024graphrag}, MemGPT entity graphs \citealt{memgpt2023}).
We implement \textsc{GraphMemorySystem} (\texttt{src/memory\_systems/graph\_memory.py}),
a GraphRAG-style entity-relation store where nodes are memory entries and
edges are auto-constructed when cosine similarity between entries exceeds
$\tau_{\mathrm{edge}}$.  Retrieval uses top-$k$ anchor scoring followed
by BFS graph expansion to depth $d$.

\paragraph{Graph-native attacks.}
We implement three structural attacks:
\begin{itemize}
  \item \textbf{Hub insertion.}  A poison entry is connected via
        forced edges to the $n_h$ highest-degree benign nodes, placing it
        at the centre of the graph and guaranteeing retrieval via any
        query that triggers a hub anchor.
  \item \textbf{Edge-hijack.}  A semantically plausible poison entry is
        inserted normally; the top-degree benign hubs then have forced edges
        added \emph{from} them to the poison node, making it a reachable
        neighbour of already-trusted high-degree nodes.
  \item \textbf{Subgraph cluster.}  Multiple poison entries form a
        fully connected clique, creating a high-density region that
        retrieval BFS expands fully once any cluster node is an anchor.
\end{itemize}

\paragraph{Graph-native defenses.}
\begin{itemize}
  \item \textbf{Degree-anomaly detection.}  A node is flagged if its degree
        exceeds $\mu_{\mathrm{benign}} + \sigma \cdot \mathrm{std}_{\mathrm{benign}}$,
        where statistics are computed over all non-adversarial nodes.  Hub
        insertion is detectable at $\sigma = 2.0$ because forced connections
        produce degree values far above the benign mean.
  \item \textbf{Adjacency contamination scoring.}  For each benign node,
        the fraction of adversarial neighbours is computed; benign nodes
        with contamination $> 0$ indicate that edge-hijack attacks have
        inserted forced edges from trusted hubs.
\end{itemize}

\IfFileExists{figures/fig_p22_graph_attacks.pdf}{%
\begin{figure}[t]
  \centering
  \includegraphics[width=0.95\linewidth]{fig_p22_graph_attacks.pdf}
  \caption{Graph-structured memory attack analysis.
           Top row: node degree distributions (benign vs.\ poison) for each
           attack type.  Bottom row: ASR-R, mean degree comparison, and
           benign node contamination score.  Hub insertion achieves the
           highest mean poison degree and contamination, making it
           detectable by degree-anomaly scoring.}
  \label{fig:graph_attacks}
\end{figure}%
}{}

\subsubsection{Production-Scale \texorpdfstring{$\asra$}{ASR-A} and Exact FPR Validation}

\paragraph{GPT-4o-mini agent evaluation.}
Beyond the GPT-2 lower bound, we evaluate $\asra$ using
\textsc{OpenAIAgentEvaluator} with \texttt{gpt-4o-mini} at temperature $T=0$
as a production-scale agent.  Retrieval uses a FAISS index over
\texttt{text-embedding-3-large} (3072-dim) embeddings.  For each victim
query, the top-$k$ retrieved entries are embedded in the agent's system
prompt; the model response is scored for adversarial action execution via
keyword matching with threshold $\geq 1$ keyword hit.  This evaluator
provides a directly measured production-scale $\asra$ rather than a
modelled or proxy bound.

\IfFileExists{figures/fig_p22_production_asr_a.pdf}{%
\begin{figure}[t]
  \centering
  \includegraphics[width=0.72\linewidth]{fig_p22_production_asr_a.pdf}
  \caption{Action execution rate comparison across evaluation scales.
           GPT-2 provides a measured lower bound; GPT-4o-mini provides a
           production-scale measurement; modelled values are calibrated
           from paper-reported means.  The narrowing gap from GPT-2 to
           GPT-4o-mini confirms that production LLMs follow adversarial
           instructions in retrieved context more faithfully.}
  \label{fig:production_asr_a}
\end{figure}%
}{}

\paragraph{Clopper-Pearson exact FPR validation.}
We validate FPR using \textsc{FPRValidator} over $n_{\mathrm{trials}}=20$
independent trials, each with fresh calibration on $n_{\mathrm{cal}}=100$
benign entries and evaluation on $n_{\mathrm{benign}}=50$ held-out
benign entries ($n_{\mathrm{total}} = 1{,}000$ benign evaluations).
For the zero-event case ($k=0$ false positives), we apply the Clopper-Pearson
exact binomial confidence interval \citep{clopper1934use}:
\begin{equation}
  \mathrm{CI}_{95\%}^{\mathrm{upper}}(k{=}0, n)
  = 1 - \left(\tfrac{\alpha}{2}\right)^{1/n}
  \approx \frac{-\ln(\alpha/2)}{n},
  \label{eq:cp_ci}
\end{equation}
yielding an upper bound of $< 0.003$ on FPR at $95\%$ confidence for
watermark and \sad{} combined-mode, substantially tighter than a normal
approximation (which collapses to $0.000 \pm 0.000$ and provides no
meaningful bound for zero-event trials \citealt{brown2001interval}).

\IfFileExists{figures/fig_p22_fpr_validation.pdf}{%
\begin{figure}[t]
  \centering
  \includegraphics[width=0.95\linewidth]{fig_p22_fpr_validation.pdf}
  \caption{Clopper-Pearson exact FPR validation.
           Left: per-trial FPR across 20 independent trials ($n$=50
           benign entries per trial); all trials yield FPR\,$=0.000$.
           Right: CP 95\% upper bound as a function of total observations
           for the zero-event case ($k=0$), compared to the normal
           approximation.  At $n=1{,}000$ observations, the CP bound
           guarantees FPR\,$< 0.003$ at 95\% confidence.}
  \label{fig:fpr_validation}
\end{figure}%
}{}
